\subsection{Funciones aleatorias continuas}
Utilizando la misma lógica de las secciones anteriores, se puede
hacer las siguientes definiciones para funciones aleatorias continuas.

Sea $P(\vec{a}) \geq 0$ una función de probabilidad de densidad
$\vec{a} \in \Omega = \set{(a_{1}, a_{2}, \ldots, a_{N}): a_i \in
\mathbb{R} }$, y denotaremos como $d^N a$ al diferencial sobre $\Omega$,
por lo que la probabilidad sobre $\Omega$ debe cumplir

\begin{gather}
  \int\dots\int_{\Omega}  P(\vec{a})\, da_1 \cdots d a_N = 1
  \label{eq:prob-aleatoria-continua}
\end{gather}

donde en vez de usar esta notación, usaremos la forma corta.

\begin{equation}
  \int_\Omega P(\vec{a})\,d^Na
  \label{eq:eq-prob-random-continue}
\end{equation}

Ahora definamos una función aleatoria continua $F: \Omega \to \mathbb{R}$
por lo que la media de $F$ será

\begin{equation}
  \bk{F} = \int_\Omega F(\vec{a}) P(\vec{a})\, d^Na
  \label{eq:mean-con-mult}
\end{equation}

El $n$-esimo momento de $F$ será

\begin{equation}
  \bk{F^n} = \int_\Omega F^n(\vec{a}) P(\vec{a}) d^N\vec{a}
  \label{eq:moment-con-mult}
\end{equation}

El $n$-simo momento central de $F$ será

\begin{equation}
  \bk{(F - \bk{F})^n}
  \label{eq:n-central-moment-con-mult}
\end{equation}

Y la desviación estandar de $F$ será

\begin{equation}
  \sigma^2 = \bk{(F - \bk{F})^2} = \bk{F^2} - \bk{F}^2
  \label{eq:std-con-mult}
\end{equation}

De manera análoga a la función delta de Kroneker se define
la función delta $\delta_\epsilon(x)$ como sigue

\begin{equation}
  \delta_\epsilon (x) =
  \begin{cases}
    0 & \text{si } x < \epsilon\\
    \frac{1}{2\epsilon} & \text{si } -\epsilon \leq x \leq \epsilon\\
    0 & \text{si } x > \epsilon
  \end{cases}
  \label{eq:delta}
\end{equation}

Por lo que la integral sobre $[-\epsilon, \epsilon] \subset
\Omega \subset \mathbb R$ será

\begin{equation}
  \int_{-\infty}^\infty \delta_\epsilon(cx - d)\, dx = \int_\Omega
  \delta_\epsilon(cx - d) dx = \frac{1}{|c|}
  \label{eq:delta-integral}
\end{equation}

Con esto la función \textit{Delta de Dirac} se puede definir como

\begin{equation}
  \delta(x) = \lim_{\epsilon \to 0} \delta_\epsilon(x)
  \label{eq:delta-dirac}
\end{equation}

y así la integral de $f(x)\delta(c(x - x_0))$ sobre $\Omega$ será

\begin{equation}
  \int_{-\infty}^\infty f(x)\delta(c(x - x_0))\,dx = \frac{f(x_0)}{|c|}
  \label{eq:dalta-dirac-integral}
\end{equation}

Asi tomando una función análitica sobre una región $x_0 \in \Omega$, y
una función $g(x)$ análitica y derivable con ceros en
$\left\{x_j\right\}_{j=1}^N$ y $g'(x_j) \neq 0 $, resulta que

\begin{equation}
  \int_{-\infty}^\infty f(x)\delta(g(x))\,dx = \sum_{j=1}^N
  \frac{f(x_j)}{|g'(x_j)|}
  \label{eq:delta-dirac-integral-2}
\end{equation}

suponiengo que se tiene dos variables aleatorias, $x$ y $y$,
junto a su densidad de probablidad, $P(x, y)$, y se define una
nueva variable aleatoria, $s$, la cual es alguna
función sobre las variables aleatorias originales, $x$ y $y$,
$s = f(x, y)$. Entonces la nueva funsión de densidad
$P(s)$ puede ser escrita como:

\begin{equation}
  P(s) = \int_{-\infty}^{\infty}\int_{-\infty}^{\infty} P(x, y)
  \delta(s - f(x, y))\, dx dy
  \label{eq:new-con-pdf}
\end{equation}

Por lo que integrando sobre $\Omega$ se obtiene la normalización

\begin{align*}
  \int_{-\infty}^{\infty} P(s)\, ds &=
  \int_{-\infty}^{\infty}\int_{-\infty}^{\infty}
  \int_{-\infty}^{\infty} P(x, y) \delta(s - f(x, y))\, dx dy ds\\
  &= \int_{-\infty}^{\infty} \int_{-\infty}^{\infty} P(x, y)\, dx dy = 1
  \label{eq:normalizacion-con-mult}
\end{align*}
